
\documentclass[12pt]{article}

\usepackage[T1]{fontenc}
\usepackage[utf8]{inputenc}
\usepackage[brazil]{babel}
\usepackage{lmodern}
\usepackage{geometry}
\usepackage{enumitem}

\geometry{margin=2.5cm}

\title{Relatório de Análise Jurídica}
\author{Sistema de IA}
\date{}

\begin{document}

\maketitle

\section*{Identificação do caso}

\textbf{Número do processo:} 0654321-09.2024.8.04.0001 \\
\textbf{UF:} AM \\
\textbf{Comarca:} Manaus \\
\textbf{Valor da causa:} R\$ 25000.0

\section*{Resumo dos fatos}
Autor aposentado e idoso alega contratação fraudulenta de empréstimo consignado via canal digital pelo Banco UFMG, sem sua anuência, com descontos indevidos em seu benefício previdenciário, requerendo declaração de inexistência do débito, restituição em dobro e indenização por danos morais.

\section*{Alegações da parte autora}
\begin{itemize}
\item Jamais contratou, solicitou ou autorizou empréstimo consignado com o Banco UFMG.
\item Não possui conta na Caixa Econômica Federal, onde foi depositado o crédito do empréstimo.
\item Nunca utilizou aplicativo ou canal digital do banco para contratação.
\item Houve uso indevido de seus dados pessoais para contratação fraudulenta.
\item Descontos indevidos comprometem sua subsistência, pois depende exclusivamente do benefício previdenciário.
\item Requereu registro de ocorrência policial e reclamação administrativa junto ao banco e Banco Central.
\end{itemize}

\section*{Documentos e evidências apresentadas pelo banco}
\begin{itemize}
\item Comprovante de crédito emitido pelo Banco UFMG confirmando operação de empréstimo consignado via aplicativo mobile.
\item Demonstrativo de evolução da dívida com parcelas pagas e saldo devedor atualizado.
\item Laudo referenciado interno do banco com evidências digitais da contratação (device fingerprint, IP, autenticação biométrica e senha eletrônica).
\end{itemize}

\section*{Pontos críticos do caso}
\subsection*{Pontos favoráveis ao banco}
\begin{itemize}
\item Apresenta comprovante de crédito e contrato formalizado com dados do autor.
\item Evidências digitais indicam contratação via aplicativo mobile com autenticação biométrica e senha.
\item Registro da operação e liberação do crédito em conta vinculada ao tomador.
\end{itemize}

\subsection*{Pontos desfavoráveis ao banco}
\begin{itemize}
\item Ausência do vídeo de liveness (biometria facial) tradicionalmente capturado no final do fluxo de contratação.
\item Conta para liberação do crédito não pertence ao autor, que nega possuir conta na Caixa Econômica Federal.
\item Autor alega não ter utilizado o aplicativo nem ter autorizado a contratação.
\item Descontos realizados em benefício previdenciário, o que agrava a situação pela vulnerabilidade do autor.
\end{itemize}

\subsection*{Observações importantes}
O caso envolve alegação de fraude com uso indevido de dados pessoais do autor idoso e aposentado, com forte indício de contratação por terceiro. A ausência do vídeo de liveness no sistema do banco pode comprometer a prova da efetiva manifestação de vontade do autor. A tutela de urgência é requerida para cessar descontos que comprometem verba alimentar.

\section*{Fatores jurídicos relevantes}

% preenchido pelo master agent
Autor pessoa idosa: sim \\
Autor aposentado/pensionista: sim \\
Pedido de tutela de urgência: sim \\
Pedido de dano moral: sim \\
Valor do dano moral: R\$ 18000.0 \\
Pedido de repetição de indébito: sim \\
Alegação de vulnerabilidade econômica: sim \\
Uso indevido de dados: sim

\section*{Análise preditiva}

% =========================
% TODO: SAÍDA DO SUBAGENT DE RISCO
% Inserir aqui:
% - probabilidade de êxito / não êxito
% - classificação preditiva
% - confiança do modelo
% =========================

\section*{Recomendação estratégica}

% =========================
% TODO: SAÍDA DO SUBAGENT DE ACORDO
% Inserir aqui:
% - recomendação de acordo ou defesa
% - valor sugerido
% - faixa de negociação
% - justificativa econômica
% =========================

\end{document}
